
\phantomsection
\addchaptertocentry{Ерөнхий дүгнэлт}
\label{Summary}

\section*{Ерөнхий дүгнэлт}
Энэхүү дипломын ажлын хүрээнд хиймэл оюунд суурилсан, шаталсан сургалтын веб системийг судалж, шинжилж, зохиомжлон, хөгжүүлэв. Судалгааны ажлын үр дүнд цахим сургалтын орчинд хэрэглэгч бүрт ижил агуулга түгээх хандлага нь програмчлалын сургалтын хувьд хангалтгүй болох нь тогтоогдож, дасан зохицох сургалт, хувь хүнд тохирсон зөвлөмж, хиймэл оюуны туслах дэмжлэг, ахицын шинжилгээ, хамтын орчин зэрэг боломжийг нэгтгэсэн платформ илүү үр дүнтэй болох нь онолын болон практик түвшинд нотлогдлоо.

Төслийн хүрээнд хэрэглэгчийн анхны түвшнийг түвшин тогтоох шалгалтаар тодорхойлох, курс болон хичээлийг түвшинтэй уялдуулан санал болгох, хичээл дуусгалт болон оноонд тулгуурлан ахиц тооцоолох, түвшин ахих шалгалтаар дараагийн шатанд ахих боломж бүрдүүлэх үндсэн логикийг хэрэгжүүлсэн. Үүний зэрэгцээ хичээлийн түвшинд хиймэл оюунаар хураангуй, сорил, дадлагын даалгавар үүсгэх, системийн түвшинд чатбот болон хяналт, хэрэглэгчийн нийт оноо, долоо хоногийн идэвх, курсын ахицыг харуулах ахицын хураангуй, сонирхол татах тоглоомын хэсэг, харилцан суралцах хамтын орчныг нэг платформд нэгтгэсэн нь уг ажлын гол үр дүн болж байна.

Техникийн хэрэгжилтийн хувьд урд талын хэсгийг React, TypeScript, Vite дээр, арын хэсгийг Django болон Django REST Framework дээр боловсруулж, өгөгдлийн сангийн түвшинд PostgreSQL болон хөгжүүлэлтийн нөөц орчноор SQLite ашиглах боломжтойгоор зохион байгуулав. Орчны тохиргоо, Docker Compose, \texttt{wait-for-db} скрипт, аюулгүй байдлын орчны хувьсагчид, JWT баталгаажуулалт, CORS/CSRF тохиргоо, тусгаарлагдсан орчны хязгаарлалт зэрэг нь энэхүү төслийг зөвхөн туршилтын интерфейсийн түвшинд бус, бодит системийн инженерчлэлийн шаардлагад ойртуулсан.

Хиймэл оюуны уялдуулалтын хувьд зөвхөн боломж нэмэх бус, нөөц логик нэмж өгснөөр системийн найдвартай байдлыг сайжруулсан нь онцлох үр дүн юм. Хураангуй, дадлагын даалгавар, хяналт зэрэг хэсгүүдэд хиймэл оюуны үйлчилгээ алдаа өгсөн ч системийн үндсэн үйл ажиллагаа тасалдахгүй байхаар шийдсэн. Мөн хамтын орчны хяналт, анхааруулгын механизм, код ажиллуулагчийн нөөцийн хязгаарлалт зэрэг нь аюулгүй байдлыг төслийн хүрээнд бодитой авч үзсэнийг харуулж байна.

Туршилт, үр дүнгийн шинжилгээнээс харахад системийн үндсэн хэрэглээний урсгалууд болох бүртгэл, нэвтрэлт, түвшин тогтоох шалгалт, курсийн шүүлт, хичээлийн дэлгэрэнгүй хуудасны харилцан үйлдэл, ахицын шинэчлэлт, хамтын орчинд нийтлэл оруулах, хяналт, хиймэл оюунаар хураангуй үүсгэх болон Docker эхлэлт зэрэг нь хоорондоо уялдаа холбоотой ажиллах боломжтой болсон. Энэ нь дасан зохицох сургалт болон хиймэл оюуны туслах сургалтын ойлголтуудыг зөвхөн онолын түвшинд бус, цогц веб системийн бодит хэрэгжилт болгон буулгаж чадсаны илрэл юм.

Цаашид энэхүү системийг хөгжүүлэхэд зөвлөмжийн хөдөлгүүрийг илүү нарийвчлах, шинжилгээний самбар нэмэх, mentor/admin-д зориулсан агуулга удирдах самбар боловсруулах, хиймэл оюунаар үүсгэсэн агуулгын чанарын үнэлгээ хийх, гар утсанд нийцсэн интерфейс болон PWA шийдэл нэвтрүүлэх, автомат тест болон CI/CD өргөтгөх, Монгол хэлний сургалтын агуулгад оновчлогдсон хиймэл оюуны ажлын урсгал бий болгох зэрэг олон боломж нээлттэй байна.

Иймд энэхүү дипломын ажил нь програмчлалын сургалтад зориулсан, хиймэл оюунд суурилсан, шаталсан сургалтын веб системийн загвар, хэрэгжилт, туршилтын суурь жишээг бүрдүүлсэн, цааш хөгжүүлж бодит хэрэглээнд ойртуулах боломжтой ажил боллоо гэж дүгнэж байна.
